\newpage

\section{Pipelined toyRISC}

The pipelined toyRISC processor design is developed in the context of a system that contains both program and data memory. Consequently, the top module contains a memory module connected to the module we are investigating, that of the pipelinedToyRISC processor.

\subsection{Structure}

\subsubsection{System level}

\includeverilogcode[caption={toyRISC-based computing system. File: toyRISCsystem.sv.  (SystemVerilog)}, 
                    label={lst:toyRISCsystem}, 
                    firstline=1, 
                    firstnumber=1, 
                    lastline=450]
                    {\verilogSourceDir/A17_toyRISCsystem.sv}

The program memory is simulated without the hardware program loading mechanism, because during the simulation process it is initialized with the user-written program in the file {\tt program.sv}, a program that is translated into binary form by the code {\tt RISCcodeGenerator.sv} which is included in the simulation program {\tt testPipelinedRISC}.

\includeverilogcode[caption={Memory system. File: memorySystem.sv.  (SystemVerilog)}, 
                    label={lst:memorySystem}, 
                    firstline=1, 
                    firstnumber=1, 
                    lastline=450]
                    {\verilogSourceDir/A17_memorySystem.sv}

The pipelined version of toyRISC represented in Figure \ref{fig46} is coded in the following module, where the three hardware structural stages (INSTRFETCH, OPFETCH, EXECUTE) of the four functional stages (INSTRFETCH, OPFETCH \& WRITEBACK, EXECUTE) of the pipeline are described. 

\includeverilogcode[caption={Pipelined toyRISC processor. File: toyRISCpipelined.sv.  (SystemVerilog)}, 
                    label={lst:toyRISCpipelined}, 
                    firstline=1, 
                    firstnumber=1, 
                    lastline=450]
                    {\verilogSourceDir/A17_toyRISCpipelined.sv}

\subsubsection{Decoder}

Decoder module works for all pipe levels providing selection signals and decoded signals. In addition to these, it also contains the automaton that manages the interrupt signal.

\includeverilogcode[caption={Decode  module in  toyRISC processor. File: 3\_DECODE.sv.  (SystemVerilog)}, 
                    label={lst:DECODE}, 
                    firstline=1, 
                    firstnumber=1, 
                    lastline=450]
                    {\verilogSourceDir/A17_DECODE.sv}

\subsubsection{First Pipe Level: Instruction Fetch}

In each clock cycle a new instruction is fetched from the program memory and inserted in the pipe. The next program address is generated according to the selection code provided by the decoder module.

\includeverilogcode[caption={Instruction fetch pipeline stage in  toyRISC processor. File: 3\_INSTRFETCH.sv.  (SystemVerilog)}, 
                    label={lst:INSTRFETCH}, 
                    firstline=1, 
                    firstnumber=1, 
                    lastline=450]
                    {\verilogSourceDir/A17_INSTRFETCH.sv}

\subsubsection{Second Pipe Level: Operands Fetch}

The second hardware level in pipe works for two functional levels in the pipe:
\begin{itemize}
    \item the second functional level of fetching the two operands from the register file, {\tt rf[0:31]}
    \item the fourth functional level of writing back the result of the operation performed by ALU in the third pipe level.
\end{itemize}

\includeverilogcode[caption={Operands fetch pipeline stage in  toyRISC processor. File: 3\_OPSFETCH.sv.  (SystemVerilog)}, 
                    label={lst:OPSFETCH}, 
                    firstline=1, 
                    firstnumber=1, 
                    lastline=450]
                    {\verilogSourceDir/A17_OPSFETCH.sv}

\paragraph{Interrupt unit} activated by the acknowledged interrupt signal, {\tt inta}. This small module of multiplexers manage to read the subroutine address and to write the return address when the processor acknowledge the interrupt signal, {\tt intIn}.

\includeverilogcode[caption={Interrupt unit in pipelined toyRISC processor. File: 4\_intUnit.sv.  (SystemVerilog)}, 
                    label={lst:intUnit}, 
                    firstline=1, 
                    firstnumber=1, 
                    lastline=450]
                    {\verilogSourceDir/A17_intUnit.sv}

\subsubsection{Third Pipe Level: Execute}

The third level of the hardware pipe executes the third level of the functional pipe which is the execution step. This hardware level contains ALU and the multiplexer which selects the right operand for ALU.

\includeverilogcode[caption={Execute pipeline stage in toyRISC processor. File: 3\_EXECUTE.sv.  (SystemVerilog)}, 
                    label={lst:EXECUTE}, 
                    firstline=1, 
                    firstnumber=1, 
                    lastline=450]
                    {\verilogSourceDir/A17_EXECUTE.sv}

\paragraph{ALU} is part of the execution level in the pipeline organization of the system.

\includeverilogcode[caption={ALU in pipelined toyRISC processor. File: 4\_ALUp.sv.  (SystemVerilog)}, 
                    label={lst:ALUp}, 
                    firstline=1, 
                    firstnumber=1, 
                    lastline=450]
                    {\verilogSourceDir/A17_ALUp.sv}

\subsection{Code generator}

For the code generator see Listing \ref{lst:RISCcodeGenerator} where the lines 377-381 must be selected and and the lines 369-373 deselected.


\subsection{Simulator}

\includeverilogcode[caption={Test pipelined toyRISC processor. File: 0\_testPipelinedRISC.sv.  (SystemVerilog)}, 
                    label={lst:testPipelinedRISC}, 
                    firstline=1, 
                    firstnumber=1, 
                    lastline=450]
                    {\verilogSourceDir/A17_testPipelinedRISC.sv}


\subsection{Testing}

For testing procedure see \ref{testProcedure}.

\includeverilogcode[caption={Programs for pipelined toyRISC processor. File: 0\_program.sv.  (SystemVerilog)}, 
                    label={lst:0_program}, 
                    firstline=1, 
                    firstnumber=1, 
                    lastline=450]
                    {\verilogSourceDir/A17_0_program.sv}



